\section{最终代码与支撑材料}
\subsection{实现入口和可复现性}
最终问题三、四入口分别为 \texttt{q3-v18}、\texttt{q4-v18}，均由 \texttt{versioned.py} 工厂构造，避免直接实例化父类时漏掉最终参数。支撑材料包含当前工作区的Python核心源码快照，原相对路径和SHA-256见 \texttt{code/MANIFEST.md}；完整程序运行仍以仓库框架和依赖为准。
\begin{longtable}{>{\raggedright\arraybackslash\small}p{4.7cm}p{9.1cm}}\toprule
主要文件&用途\\\midrule\endhead
\texttt{geometry.py}&楔形裁剪、圆盘外包络、区域直径及最小包围圆。\\
\texttt{knowledge.py}&频道—路径稀疏矩阵、真实坐标和反馈摘要。\\
\texttt{knowledge\_layer.py}&失败清除排除圆、相交方格、覆盖计划及代价比较。\\
\texttt{strategy\_v15.py}&问题三的负观测估计修正、2-opt与or-opt调度。\\
\texttt{strategy\_v17.py}&问题四24点布局，继承v14射线搜索和在线调度。\\
\texttt{strategy\_v18.py}&两问最终策略，共用知识矩阵和覆盖清除混入模块。\\
\texttt{versioned.py}&交付入口与参数：问题三6/1130/3、or-opt四轮。\\
{\scriptsize\texttt{q3\_or\_opt\_ablation.py}}&固定计时和种子，分离布局、路由和覆盖清除的收益。\\
{\tiny\texttt{verify\_layout\_certificate.py}}&连续单元覆盖证书和有限采样交叉检验。\\
\texttt{make\_assets.py}&按实际数据生成论文SVG/PDF图和实验表。\\\bottomrule
\end{longtable}

\subsection{离线复现实验命令}
在项目根目录执行下列命令；这些命令只访问本地模拟器和文件。
\begin{Verbatim}[fontsize=\scriptsize,frame=single]
python3 script/run.py --strategy q3-v18 --seed 1
python3 script/run.py --strategy q4-v18 --directional --seed 1
python3 script/q3_or_opt_ablation.py --json paper/build/q3-final-ablation.json
python3 script/verify_layout_certificate.py --layout v17
python3 paper/make_assets.py
python3 paper/verify_formulas.py
bash paper/build.sh
\end{Verbatim}
第三问新复跑结果保存在 \texttt{paper/data/q3-final-ablation.json}；第四问已有逐例记录位于 \texttt{logs/audit/}，论文保留其汇总与源文件哈希。演练汇总保存在 \texttt{paper/data/official-drill-q18.json}，不能代替要求提交的正式原名日志。

\subsection{覆盖式清除核心代码摘录}
以下三个函数省略注释与文档字符串，执行语句直接取自当前 \path{knowledge_layer.py}，辅助几何、类型定义和常量见完整快照。相交格不等于仅格心落入区域，整格剔除也不能简化为格心在失败圆内。覆盖保证还要求保守域包含真值、无不可达格和计划实际执行完成。
\lstinputlisting[language=Python]{code/excerpt_cover_clear.py}
